\documentclass[border=2pt]{standalone}
% Estilo comum dos fluxogramas do Caderno C3 (cores do modelo do MOP)
\usepackage[T1]{fontenc}
\usepackage[utf8]{inputenc}
\usepackage{helvet}
\renewcommand{\familydefault}{\sfdefault}
\usepackage{amssymb}
\usepackage{tikz}
\usetikzlibrary{positioning, arrows.meta, shapes.geometric, shapes.symbols, shapes.misc,
  calc, fit, backgrounds, decorations.pathmorphing, patterns, shadows}
% Paleta do modelo do caderno
\definecolor{cTeal}{HTML}{00747A}
\definecolor{cNavy}{HTML}{0B3A5A}
\definecolor{cBlue}{HTML}{0A73BA}
\definecolor{cBlueBg}{HTML}{EBF4F9}
\definecolor{cTealBg}{HTML}{E3F1F2}
\definecolor{cTealMid}{HTML}{9CCFD2}
\definecolor{cGrayBg}{HTML}{F2F2F2}
\definecolor{cGrayLine}{HTML}{BFBFBF}
\definecolor{cGrayText}{HTML}{595959}
\definecolor{cText}{HTML}{1A1A1A}
\definecolor{cAmber}{HTML}{B7791F}
\definecolor{cAmberBg}{HTML}{FDF3DC}
\definecolor{cGreen}{HTML}{3D7D3A}
\definecolor{cGreenBg}{HTML}{E8F2E4}
\color{cText}
\tikzset{
  every node/.style={font=\footnotesize, text=cText},
  etapa/.style={draw=cTeal, fill=cTealBg, rounded corners=2pt, align=center,
     inner sep=3.5pt, minimum height=0.85cm, text width=#1, line width=0.6pt},
  etapa/.default={3.2cm},
  marco/.style={draw=cNavy, fill=cNavy, rounded corners=2pt, align=center,
     inner sep=3.5pt, minimum height=0.85cm, text width=#1,
     font=\footnotesize\bfseries, text=white},
  marco/.default={3.2cm},
  marcoteal/.style={draw=cTeal, fill=cTeal, rounded corners=2pt, align=center,
     inner sep=3.5pt, minimum height=0.85cm, text width=#1,
     font=\footnotesize\bfseries, text=white},
  marcoteal/.default={3.2cm},
  externo/.style={draw=cGrayLine, fill=cGrayBg, rounded corners=2pt, align=center,
     inner sep=3.5pt, minimum height=0.85cm, text width=#1},
  externo/.default={3.2cm},
  produto/.style={draw=cTeal!70, fill=white, rounded corners=1pt, align=center,
     inner sep=2.5pt, minimum height=0.6cm, text width=#1, font=\scriptsize},
  produto/.default={2.4cm},
  decisao/.style={draw=cAmber, fill=cAmberBg, diamond, aspect=2, align=center,
     inner sep=1pt, font=\scriptsize, line width=0.6pt},
  correcao/.style={draw=cAmber, fill=cAmberBg, rounded corners=2pt, align=center,
     inner sep=3pt, minimum height=0.8cm, text width=#1, font=\scriptsize},
  correcao/.default={2.6cm},
  transversal/.style={draw=cBlue, fill=cBlueBg, dashed, rounded corners=2pt, align=center,
     inner sep=3.5pt, text width=#1, line width=0.6pt},
  transversal/.default={6cm},
  seta/.style={-{Stealth[length=2.2mm,width=1.8mm]}, line width=0.8pt, draw=cNavy},
  setaleve/.style={-{Stealth[length=1.8mm,width=1.5mm]}, line width=0.6pt, draw=cGrayText},
  setatracejada/.style={-{Stealth[length=2mm]}, line width=0.7pt, dashed, draw=cGrayText},
  rotulo/.style={font=\scriptsize\itshape, text=cGrayText, fill=white, inner sep=1pt},
  fase/.style={signal, signal to=east, signal from=west, draw=cTeal, fill=cTealBg,
     align=center, font=\scriptsize\bfseries, minimum height=1.05cm, text width=#1,
     inner sep=2pt},
  fase/.default={1.9cm},
  raia/.style={draw=cGrayLine, fill=cGrayBg, minimum width=#1, minimum height=0.9cm,
     anchor=north west},
  tituloraia/.style={font=\scriptsize\bfseries, text=cNavy, align=center, text width=1.75cm},
  cab/.style={font=\scriptsize\bfseries, text=white, fill=cNavy, align=center,
     inner sep=3pt, text width=#1, minimum height=0.6cm},
  cab/.default={3cm},
  bloco/.style={draw=cGrayLine, fill=cGrayBg, rounded corners=3pt},
}
% sem hifenização nas caixas dos fluxogramas
\AtBeginDocument{\hyphenpenalty=10000 \exhyphenpenalty=10000 \sloppy}

\begin{document}

\begin{tikzpicture}[x=1cm,y=1cm,
   num/.style={circle, fill=petroleo, text=white, font=\fonte{8}{9}\bfseries, inner sep=0pt,
     minimum size=0.44cm},
   traco/.style={draw=petroleo, line width=0.8pt},
   cano/.style={draw=cinza60, line width=2pt},
   nota/.style={font=\fonte{7.5}{8.8}, text=cinza80}]
  \fill[cinza5] (0,0) rectangle (15.6,-2.9);
  \draw[cinza60, line width=1.2pt] (0,0) -- (15.6,0);
  % domicílio
  \filldraw[traco, fill=white] (0.3,0) rectangle (2.1,1.3);
  \filldraw[traco, fill=white] (0.1,1.3) -- (1.2,2.0) -- (2.3,1.3) -- cycle;
  \draw[traco, line width=0.6pt] (0.95,0) rectangle (1.45,0.75);
  \node[num] at (1.2,2.45) {1};
  % tubulação
  \draw[cano] (1.6,-0.45) -- (2.95,-0.45);
  % caixa de inspeção
  \fill[white] (2.95,-0.15) rectangle (3.55,-0.85);
  \draw[traco] (2.95,0) -- (2.95,-0.85) -- (3.55,-0.85) -- (3.55,0);
  \node[num] at (3.25,0.5) {2};
  \draw[cano] (3.55,-0.45) -- (4.3,-0.45);
  % tanque séptico com duas câmaras
  \fill[white] (4.3,-0.1) rectangle (7.3,-2.0);
  \fill[petroleo12] (4.32,-0.6) rectangle (7.28,-1.6);
  \fill[pattern=north east lines, pattern color=petroleo!70] (4.32,-1.6) rectangle (7.28,-1.98);
  \draw[traco] (4.3,0) -- (4.3,-2.0) -- (7.3,-2.0) -- (7.3,0);
  \draw[traco, line width=0.7pt] (5.9,-0.1) -- (5.9,-1.5);
  \node[nota, fill=white, inner sep=1pt] at (6.6,-1.79) {lodo};
  \node[num] at (5.8,0.5) {3};
  \draw[cano] (7.3,-0.5) -- (8.05,-0.5);
  % filtro anaeróbio
  \fill[white] (8.05,-0.1) rectangle (10.05,-2.0);
  \fill[pattern=crosshatch dots, pattern color=petroleo] (8.07,-0.8) rectangle (10.03,-1.75);
  \draw[traco] (8.05,0) -- (8.05,-2.0) -- (10.05,-2.0) -- (10.05,0);
  \node[num] at (9.05,0.5) {4};
  \draw[cano] (10.05,-0.55) -- (11.0,-0.55);
  % disposição final no solo
  \draw[traco, line width=0.6pt] (11.0,-0.3) -- (11.0,-1.0) -- (15.1,-1.0) -- (15.1,-0.3);
  \draw[cano, dash pattern=on 4pt off 2pt] (11.0,-0.55) -- (15.0,-0.55);
  \foreach \x in {11.5,12.3,13.1,13.9,14.7}
    \draw[-{Latex[length=1.6mm]}, draw=petroleo, line width=0.7pt] (\x,-1.05) -- (\x,-1.85);
  \node[num] at (13.05,0.5) {5};
  % legenda
  \begin{scope}[shift={(0.1,-3.35)}]
    \foreach \n/\t/\x/\y in {1/Domicílio/0/0, 2/Caixa de inspeção/3.2/0,
        3/{Tanque séptico, com remoção periódica do lodo}/7.4/0, 4/Filtro anaeróbio/0/-0.5,
        5/{Disposição final no solo (vala de infiltração ou sumidouro)}/3.2/-0.5} {
      \node[num, minimum size=0.36cm, font=\fonte{7}{8}\bfseries] at (\x+0.18,\y) {\n};
      \node[anchor=west, font=\fonte{8}{9.4}] at (\x+0.42,\y) {\t};
    }
  \end{scope}
  \node[text=cinza80, anchor=north, text width=15cm, align=center] at (7.8,-4.25)
    {Exemplo ilustrativo. A solução de cada domicílio é definida no projeto individual, conforme a
     ABNT NBR 17076:2024, as condições do terreno e o número de moradores.};
\end{tikzpicture}

\end{document}
